\uuid{YeNv}
\titre{Primitives usuelles}
\auteur{hanani ; mzari}
\organisation{exo7}
\licence{CC BY-NC-SA 4.0}
\datecreate{2020-01-20}
\niveau{L1}
\module{Analyse}
\chapitre{Calcul d'intégrales}
\sousChapitre{Primitives diverses}
\difficulte{1}
\theme{Primitives, primitives usuelles}
\exo7qcmid{lille-2.primitives.008}

\contenu{
\begin{question}
\qclassification{127.11}
\qtags{motcle=primitives usuelles}
Parmi les assertions suivantes, cocher celles qui sont vraies :
\begin{answers}
\bad{Une primitive de $\displaystyle \cos (\pi x)+\frac{1}{x}$ sur $]0,+\infty[$ est $\displaystyle \sin (\pi x)+\ln x$.}
\good{Une primitive de $\displaystyle \cos (\pi x)+\frac{1}{x}$ sur $]0,+\infty[$ est $\displaystyle \frac{\sin (\pi x)}{\pi}+\ln (\pi x)$.}
\bad{La primitive de $\displaystyle \cos (\pi x)+\frac{1}{x}$ sur $]0,+\infty[$ qui s'annule en $1$ est $\displaystyle \sin (\pi x)+\ln x$.}
\good{La primitive de $\displaystyle \cos (\pi x)+\frac{1}{x}$ sur $]0,+\infty[$ qui s'annule en $1$ est $\displaystyle \frac{\sin (\pi x)}{\pi}+\ln x$.}
\end{answers}
\begin{explanations}
Par linéarité, les primitives de $\displaystyle \cos (\pi x)+\frac{1}{x}$ sur $]0,+\infty[$ sont
$$F(x)=\int \left(\cos (\pi x)+\frac{1}{x}\right)\mathrm{d}x=\frac{\sin (\pi x)}{\pi}+\ln x+k,\; k\in \Rr.$$
Ensuite $F(1)=0\Rightarrow k=0$. Donc la primitive de $\displaystyle \cos (\pi x)+\frac{1}{x}$ sur $]0,+\infty[$ qui s'annule en $1$ est $\displaystyle F(x)=\frac{\sin (\pi x)}{\pi}+\ln x$.
\end{explanations}
\end{question}
}
